\documentclass[10pt,a4paper]{amsart}
\usepackage[margin=19mm]{geometry}
\usepackage{amsmath,amssymb,mathtools}
\usepackage[scr=rsfs,cal=euler]{mathalfa}
\usepackage{xfrac}
\usepackage[unicode,hidelinks,bookmarks=false]{hyperref}
\newcommand{\ie}{i.e.\ }
\newcommand{\cf}{cf.\ }
\newcommand{\resp}{respectively\ }
\newcommand{\Subsection}{\subsection}
\providecommand{\nobreakdash}{\nobreak\hbox{-}}
\setlength{\emergencystretch}{2em}
\multlinegap0pt
\usepackage{amssymb}
\usepackage{xcolor}
\usepackage{cancel}
\newtheorem{lemma}{Lemma}
\newtheorem{theorem}{Theorem}
\newcommand{\N}{\mathbb{N}}
\newcommand{\R}{\mathbb{R}}
\newcommand{\Z}{\mathbb{Z}}
\newcommand{\dis}{\operatorname{dis}}
\newcommand{\nets}{\operatorname{Net}}
\renewcommand{\ss}{\subset}
\title{Contractibility of the space of $\varepsilon$-nets in $\R$}
\author{Ivan N. Mikhailov}
\date{Source: arXiv:2605.19680v1 (CC BY 4.0)}
\begin{document}
\maketitle

\noindent\emph{Source and licence.} Contractibility of the space of $\varepsilon$-nets in $\R$. Version arXiv:2605.19680v1 (CC BY 4.0), distributed by arXiv under the Creative Commons Attribution 4.0 International licence, http://creativecommons.org/licenses/by/4.0/, as recorded in the arXiv licence metadata for this exact version and shown on the abstract page https://arxiv.org/abs/2605.19680v1. Full licence notice: \texttt{licenses/arxiv-2605.19680-CC-BY-4.0.txt}.\par\smallskip

\noindent\textit{Editorial scope.} Counted unit: the article's theorem theorem (source0.tex lines 305--307). The article's complete original proof of it is delivered with it (source0.tex lines 309--379, one \texttt{proof} environment of 71 lines). No mathematics is written, repaired or re-derived: the statement and the proof are the article's own lines, in the article's own notation, and no retained line is substituted or re-flowed.  Retained uncounted as the source-local context the proof needs: lines 289--291.  Nothing was omitted from any retained span, so the package carries no omission record.  No retained line contains a citation, so the wrapper carries no bibliography and no bibliography entry.  No retained line contains a figure environment, an image-inclusion command or drawing code, so no image is delivered and the package declares no asset dependency.  Editorial formatting only, in addition: an AMS \texttt{amsart} preamble that restates the article's own declarations and macros used by the retained lines, namely the environments \texttt{lemma}, \texttt{theorem} on separate counters, together with the standard packages those lines need.  The retained environments print in this document as Lemma 1, Theorem 1 in order of appearance, whereas the article prints the same items with the numbering of its own full document.  The complete native source of the article as distributed in the arXiv e-print for this exact version is bundled unchanged as \texttt{sources/200934/source0.tex}.
\begin{lemma}\label{lemma: transfer_separated}
Suppose $R\in\mathcal{R}(A, B)$, for some subsets $A,B\R$, $\dis R = c$, and $A$ is $t$-separated, for some $t > 2c$. For each $p\in A$, we choose $p'\in R(p)$ arbitrarily. Then, for every three distinct points $p, q, r\in \R$, if $p$ lies between $q$ and~$r$, then $p'$ lies between $q'$ and~$r'$.
\end{lemma}

\begin{theorem}
The space $\nets_{GH}(\R)$ is contractible.
\end{theorem}

\begin{proof}
Define $\Phi\: \nets(\R)\times [0, 1]\to \nets(\R)$ as follows: $$\Phi(X, \lambda) = \begin{cases}B_{\lambda/(1-\lambda)}(X) & \lambda\in[0, 1),\\ \R &\lambda = 1.\end{cases}$$ 
% In other words, $\Phi(X, \lambda) = \overline{\cup_{x\in X}[x-\lambda, x+\lambda]}$. 

\begin{lemma}\label{lemma:contraction}
A mapping $\Phi$ is continuous. 
\end{lemma}

\begin{proof}
It suffices to show that $\Phi$ is continuous with respect to each of two variables $X$ and $\lambda$. Once again, denot $f(x) = \frac{x}{1-x}$.

1) Take arbitrary $\{\lambda_n\}_{n\in\N}$ such that $\displaystyle \lim_{n\to\infty}\lambda_n = \lambda$, where $\lambda\in[0, 1)$. Then $\displaystyle \lim_{n\to\infty}f(\lambda_n) = f(\lambda)$. Consider a natural correspondence $R$ between $B_{f(\lambda_n)}(X)$ and $B_{f(\lambda)}(X)$ which is the union of natural correspondences between segments $[x- f(\lambda_n), x+f(\lambda_n)]$ and $[x-f(\lambda), x + f(\lambda)]$. Note that 
\begin{align*}
\dis R\le 2|f(\lambda_n) - f(\lambda)| = 2\Big|\frac{\lambda_n - \lambda}{(1-\lambda_n)(1-\lambda)}\Big|. 
\end{align*}
 
Thus, $$\lim_{n\to\infty}d_{GH}\Bigl(\Phi(X,\lambda_n), \Phi(X, \lambda)\Bigr) = 0.$$

It remains to show that, for a sequence $\lambda_n$ tending to~$1$, we have $\displaystyle\lim_{n\to\infty} d_{GH}(\Phi(\lambda_n, X), \R) = 0$. Since $X$ is an $\varepsilon$-net in $\R$, we have $d_H(X, \R)<\infty$. Thus, there exists such $t$ that, for any $t' > t$, we have $B_{t'}(X) = \R$, and, hence, the desired statement is true.

2) Take arbitrary $\{X_n\}_{n\in\N}$ such that $\displaystyle \lim_{n\to\infty}X_n = X$. Let us show that $\displaystyle \lim_{n\to\infty}\Phi(X_n, \lambda) = \Phi(X,\lambda)$.

Choose arbitrary $R_n\in\mathcal{R}(X, X_n)$. By passing to a subsequence, we may assume that $\dis R_n \le \frac{2}{n+100}\lambda$. In particular, $d_{GH}(X_n, X)\le \frac{1}{n}$.

Fix some $100$-separated infinite sequence of points $X' = \{p_n\}_{n\in\Z}\ss X$ such that $p_n < p_m \Longleftrightarrow n < m$. Choose $p_k'\in R_n(p_k)$. By Lemma~\ref{lemma: transfer_separated}, the order of points $p_k'$ on the real line is the same (or inverted) as the order of points $p_k$ on the real line. In the remaining proof we assume that these orders coincide with the usual order ($<$) on the real line.

\begin{lemma}\label{lemma: order_change} Suppose $(a, a')\in R_n$, $(b, b')\in R_n$ satisfy inequalities $a < b$, $a' > b'$. Then $b-a \le 2\dis R_n$. 
\end{lemma}

\begin{proof}
Choose $p_k$ such that $p_k > b+100$. Then, by definition of distortion, $\big||p_na|-|p_n'a'|\big|\le \dis R_n$, $\big||p_nb|-|p_n'b'|\big|\le \dis R_n$. Hence, $$0 > b' - a' = |p'-b'| - |p'-a'|\ge |p - b| - |p-a| - 2\dis R_n = b-a - 2\dis R_n \Longleftrightarrow b-a < 2\dis R_n.$$
\end{proof}

Now we extend $R_n$ to a correspondence $R_n'$ between $B_{f(\lambda)}(X)$ and $B_{f(\lambda)}(X_n)$ such that $\dis R_n'\le 5\dis R_n$. 

Take arbitrary $a\in B_{f(\lambda)}(x)$ for some $x\in X$. Unless $a\in\pi_X(R_n)$, we take arbitrary $x'\in R_n(x)$ and add $(a, x' + a - x)$ to $R_n$. 

Let us check the stated inequality. 

We consider three cases:

1) $(a, x' + a-x)$ and $(b, y'+b -y)$, where $a\in B_{f(\lambda)}(x)$, $b\in B_{f(\lambda)}(y)$, $(x, x'), (y, y')\in R_n$.

Here we consider two subcases. 

Wlog, we assume that $x \le y$. 
% (in case $x = y$ everything is clear?). 

\textbf{Case 1.1}. $x' \le y'$. Then 
\begin{multline*}
\big||a-b| - |x'+a-x - y'-b + y|\big| = \big||x + a-x - y - b + y| - |x'+a-x - y'-b + y|\big| \le\\\le \big||x-y| - |x'-y'|\big|\le \dis R_n.
\end{multline*}

\textbf{Case 1.2}. $x' \ge y'$. By Lemma~\ref{lemma: order_change}, we know that $y-x<2\dis R_n < \frac{\lambda}{4}$. Thus,
\begin{align*}
\big||a-b| - |x'+a-x - y'-b + y|\big| \le \big||a-b| - |a-b|\big| + |x-y|+|x'-y'|\le 2|x-y| + \dis R_n\le  5\dis R_n.
\end{align*}


2) $(a, x' + a-x)$ and $(r, r')$, where $a\in B_{f(\lambda)}(x)$, $(x, x'), (r, r')\in R_n$.

This case can be considered similarly to the first one.

3) $(r, r')$ and $(s, s')$ where $(r, r'), (s, s')\in R_n$. 

In this case $\big||r-s|-|r'-s'|\bigr|\le \dis R_n$, by definition of distortion.
\end{proof}


From Lemma~\ref{lemma:contraction}, it follows that $\Phi$ is the desired retraction of $\nets_{GH}(\R)$ onto $\R$.
\end{proof}

\end{document}
